% Supplemental Material fragment for the Muru Comment: the correlation table
% and its accompanying text. Numbers emitted programmatically from
% stage4_stats.json plus the shape-only statistic (see
% results/muru_quantitative_statements_2026-08-31.md); do not hand-edit the
% values. 216 views = 36 azimuths x 6 halos, FWHM 2 deg unless stated.
% 2026-09-16: gNFW^2 row regenerated from the CELL-AVERAGED template
% (stage4_stats.py rerun; analysis/docs/muru_gnfw2_quadrature_2026-09-08.md).
% Only that row moved: J 0.92->0.91, stars 0.72->0.70 at 2 deg. The 3-deg
% ordering (gNFW^2 0.957 vs Coleman 0.949) survives rounding.
% Numbered \subsection so it sits as A.3 between the concentration
% measurements and the numerical checks.

\subsection{Map correlations, and why the comparison does not rest on them}

A natural summary of morphological similarity is the Pearson correlation
between each map and each template; Table~\ref{tab:corr} gives it. We report
it because it does \emph{not} discriminate, and the reason is instructive.

On the stellar side the correlation behaves as expected: the old-stellar maps
prefer each of the three bulge templates over the spherical gNFW$^{2}$
profile in all 216 views, with median correlations of $0.94$--$0.98$ against
$0.70$. On the annihilation side it does not: the corrected maps correlate
with every template at $0.91$--$0.93$, the Coleman template ($0.93$) lying
above gNFW$^{2}$ ($0.91$). All of these maps are centrally peaked,
and a correlation between linear maps is dominated by that shared structure
rather than by the differences of shape and radial profile at issue. At
$3^{\circ}$ smoothing the ordering reverses again (gNFW$^{2}$ $0.96$,
Coleman $0.95$), confirming that the statistic is not measuring a stable
property of the maps.

Dividing each map by its own azimuthally averaged radial profile isolates the
angular structure and makes the degeneracy explicit. The corrected
annihilation maps then correlate with the Cao template at $0.92$, essentially
as well as the old-stellar maps do ($0.95$). This is not evidence that the two
morphologies agree; it reflects the fact stated in the text that the corrected
maps remain flattened toward the disk plane in the same sense as the bulge.
The distinction between them is one of degree and of radial profile, and
neither is captured by a correlation coefficient. (The gNFW$^{2}$ entries in
those columns are not meaningful: a spherical template has no angular
structure to correlate against.)

The comparisons in the text therefore use the concentration $C$ and the
isophote axis ratios, which are not degenerate in this way.

% table* (full width): with two range columns the table is 37 pt too wide
% for a single column (18 pt even at \footnotesize) -- measured 2026-09-16.
\begin{table*}[b]
\caption{Median Pearson correlation of the corrected annihilation maps
($J$) and of the old-stellar maps with each template, over the 216
viewing geometries at $2^{\circ}$ smoothing; ranges in brackets.
``Angular structure'' divides each map by its own azimuthally
averaged radial profile over $0.5^{\circ}$--$15^{\circ}$.}
\label{tab:corr}
\begin{ruledtabular}
\begin{tabular}{lcccc}
 & \multicolumn{2}{c}{full map} & \multicolumn{2}{c}{angular structure} \\
\cline{2-3}\cline{4-5}
template & $J$ & stars & $J$ & stars \\ \hline
gNFW$^{2}$ ($\gamma=1.2$) & 0.91 [0.84, 0.95] & 0.70 [0.16, 0.85] & 0.11 & 0.04 \\
Freudenreich & 0.91 [0.86, 0.95] & 0.94 [0.22, 0.97] & 0.65 & 0.64 \\
Cao \textit{et al.} & 0.92 [0.86, 0.97] & 0.98 [0.23, 0.99] & 0.92 & 0.95 \\
Coleman \textit{et al.} & 0.93 [0.88, 0.97] & 0.96 [0.22, 0.98] & 0.69 & 0.71 \\
\end{tabular}
\end{ruledtabular}
\end{table*}
